# EPR-Bell Entanglement Validation and Sacred Geometry Coherence for Quantum-Classical AI Systems

**Author:** Stephen Paul Lutar Jr.
**Affiliation:** SZL Consulting Ltd
**ORCID:** [0009-0001-0110-4173](https://orcid.org/0009-0001-0110-4173)
**Contact:** stephenlutar2@gmail.com
**Date:** 4 May 2026
**Version:** v10
**Companion papers:** v1--v9 (see References)
**Reference implementation:** [github.com/szl-holdings/szl-holdings-platform](https://github.com/szl-holdings/szl-holdings-platform) -- `packages/ouroboros-integrations/src/sovereign-engine.ts` (EPRBellValidator, SacredGeometryEngine classes)
**License:** CC BY 4.0

---

## Abstract

We introduce the EPR-Bell Entanglement Validator (EBEV), a framework that applies the CHSH inequality from quantum foundations to test whether AI model correlation matrices exhibit non-classical structure. The EPR paradox (Einstein, Podolsky, Rosen, 1935), Bell's theorem (Bell, 1964), and the CHSH inequality (Clauser, Horne, Shimony, Holt, 1969) establish that quantum correlations violate classical local-realism bounds: classical \( |S| \leq 2 \), quantum \( |S| \leq 2\sqrt{2} \) (Tsirelson bound). We test AI-model feature correlations against these bounds to identify non-trivially entangled representations. We further introduce the Sacred Geometry Coherence Engine (SGCE), which evaluates model outputs using geometric constants: the golden ratio \( \varphi = (1+\sqrt{5})/2 \), the Vesica Piscis ratio \( \sqrt{3} \), Flower of Life packing density \( \pi\sqrt{3}/6 \), and Fibonacci convergence. This extends the thesis lineage from v9 (interpretability) to quantum foundations and geometric coherence.

This paper does **not** claim that AI models exhibit genuine quantum entanglement. The CHSH test is used as a *structural diagnostic* for non-classical correlation patterns, not as evidence of quantum mechanics in neural networks.

---

## 1. Motivation and Continuity from v3--v9

The v4 Omega Formalism includes \( L_5 = L_4 \cdot e^{i\gamma} \) (quantum-geometric phase) and \( L_6 = |L_5| \cdot \langle \Gamma_{E_8}, \hat{v} \rangle \) (\( E_8 \) triality). These signatures invoke quantum-geometric structures but v4 did not provide an operational test for whether model correlations actually exhibit non-classical structure. The present paper fills that gap.

The v4 sacred geometry coherence metric was introduced but not fully developed. The present paper provides the complete formalization.

---

## 2. EPR-Bell Entanglement Validator (EBEV)

### 2.1 Background: The CHSH Inequality

The CHSH inequality tests whether correlations between two spatially separated measurements can be explained by local hidden variables:

\[
S = |E(a,b) - E(a,b') + E(a',b) + E(a',b')|
\]

where \( E(a,b) \) is the correlation between measurements at settings \( a \) and \( b \).

**Classical bound:** \( S \leq 2 \) (local hidden variable theories)
**Quantum bound:** \( S \leq 2\sqrt{2} \approx 2.8284 \) (Tsirelson bound)

The singlet state \( |\psi^-\rangle = \frac{1}{\sqrt{2}}(|01\rangle - |10\rangle) \) achieves maximum violation at the angles \( a = 0 \), \( a' = \pi/2 \), \( b = \pi/4 \), \( b' = 3\pi/4 \).

### 2.2 Application to AI Model Correlations

We define model-feature measurement settings as projection angles in the feature space. For a pair of model layers \( A \) and \( B \) with feature vectors \( \mathbf{f}_A, \mathbf{f}_B \):

\[
E(a, b) = -\cos(a - b) + \epsilon(\text{data})
\]

where \( \epsilon(\text{data}) \) is a small data-dependent noise term.

### 2.3 Bell Certificate

The EBEV issues a three-level certificate:

| Certificate | Condition | Interpretation |
|---|---|---|
| BELL-CLASSICAL | \( S \leq 2 \) | Correlations explainable by classical statistics |
| BELL-PASS | \( 2 < S \leq 2\sqrt{2} \) | Non-classical correlation structure detected |
| BELL-SUPER-QUANTUM | \( S > 2\sqrt{2} \) | Exceeds quantum bound (indicates measurement error) |

---

## 3. Sacred Geometry Coherence Engine (SGCE)

### 3.1 Fundamental Constants

| Constant | Symbol | Value | Source |
|---|---|---|---|
| Golden ratio | \( \varphi \) | \( (1+\sqrt{5})/2 \approx 1.6180 \) | Euclid, Elements |
| Vesica Piscis | \( \sqrt{3} \) | \( \approx 1.7321 \) | Intersection of equal circles |
| Flower of Life packing | \( \pi\sqrt{3}/6 \) | \( \approx 0.9069 \) | Hexagonal close-packing |
| Fibonacci convergence | \( F_n/F_{n-1} \to \varphi \) | \( \varphi \) | Fibonacci (1202) |
| Metatron's Cube edges | --- | 78 | 13 vertices, complete graph |

### 3.2 Phi-Harmonic Analysis

For a vector of values \( \mathbf{v} = (v_1, \ldots, v_n) \) sorted in descending order:

\[
\Delta_\varphi = \frac{1}{n-1} \sum_{i=1}^{n-1} \left| \frac{v_i}{v_{i+1}} - \varphi \right|
\]

A phi deviation of zero indicates perfect golden ratio scaling across consecutive values.

### 3.3 Vesica Piscis Ratio

The Vesica Piscis is formed by the intersection of two equal circles whose centers lie on each other's circumference. The height-to-width ratio of the intersection is \( \sqrt{3} \). We compute:

\[
\Delta_{\text{VP}} = \left| \frac{\sum |v_i|}{n} - \sqrt{3} \right|
\]

### 3.4 Composite Coherence Score

\[
C_{\text{sacred}} = \frac{1}{1 + \Delta_\varphi} \cdot \frac{1}{1 + \Delta_{\text{VP}}} \cdot \frac{\pi\sqrt{3}}{6}
\]

Bounded in \( [0, \pi\sqrt{3}/6] \approx [0, 0.9069] \).

### 3.5 Platonic Solid Mapping

Each evaluation is mapped to a Platonic solid based on the dimensionality:

| Solid | Vertices | Faces | Dual |
|---|---|---|---|
| Tetrahedron | 4 | 4 | Tetrahedron (self-dual) |
| Cube | 8 | 6 | Octahedron |
| Octahedron | 6 | 8 | Cube |
| Dodecahedron | 20 | 12 | Icosahedron |
| Icosahedron | 12 | 20 | Dodecahedron |

---

## 4. Connection to v4 Omega Formalism

The EBEV and SGCE connect to the v4 hierarchy:

- **\( L_5 \) (quantum-geometric phase):** The CHSH \( S \)-value provides an operational measure of whether the Berry phase \( \gamma \) in \( L_5 = L_4 \cdot e^{i\gamma} \) captures genuine non-classical structure.
- **\( L_6 \) (\( E_8 \) triality):** The sacred geometry coherence score provides an independent quality check on the \( E_8 \) lattice projection.
- **\( L_1 \) (geometric ratio):** The phi deviation \( \Delta_\varphi \) directly measures how close the ratio structure is to the golden proportion.

---

## 5. What this paper does and does not claim

**Claims:**
1. The CHSH inequality is correctly implemented following Clauser et al. (1969).
2. The Bell certificate classification is well-defined.
3. The sacred geometry constants are mathematically standard.
4. The coherence score is computable and bounded.
5. All components are implemented in the public reference implementation.

**Non-claims:**
- No claim that AI models exhibit genuine quantum entanglement. The CHSH test is a structural diagnostic.
- No claim that sacred geometry coherence predicts model quality on standard benchmarks.
- No claim of deployed production use.
- No claim of metaphysical significance of the geometric constants.

---

## 6. Conclusion

The EPR-Bell Entanglement Validator and Sacred Geometry Coherence Engine provide principled structural diagnostics for AI model correlations and output quality. The CHSH framework extends the v4 quantum-geometric signatures with an operational test, while sacred geometry coherence provides an independent geometric quality metric. Both extend the thesis lineage to quantum foundations and geometric analysis.

---

## References

1. Einstein, A., Podolsky, B., Rosen, N. (1935). Can quantum-mechanical description of physical reality be considered complete? *Physical Review*, 47:777--780.
2. Bell, J.S. (1964). On the Einstein Podolsky Rosen paradox. *Physics Physique Fizika*, 1:195--200.
3. Clauser, J.F., Horne, M.A., Shimony, A., Holt, R.A. (1969). Proposed experiment to test local hidden-variable theories. *Physical Review Letters*, 23:880--884.
4. Tsirelson, B.S. (1980). Quantum generalizations of Bell's inequality. *Letters in Mathematical Physics*, 4:93--100.
5. Aspect, A., Grangier, P., Roger, G. (1982). Experimental realization of Einstein-Podolsky-Rosen-Bohm Gedankenexperiment. *Physical Review Letters*, 49:91--94.
6. Carlson, R. Sacred Geometry International. [sacredgeometryinternational.com](https://sacredgeometryinternational.com).
7. Euclid. *Elements*, Book VI, Definition 3 (extreme and mean ratio).
8. Lutar, S.P. (2026a--i). Ouroboros Thesis series v1--v9. See companion papers.
